\begin{table}[t]
\centering
\small
\caption{Capabilities of each method, independent of embedding quality.
\emph{Parametric} means unseen points can be projected onto an existing map.
\emph{Landscape grid}: the prior places $K$ components at fixed, known
positions, so any quantity---measured activity, predicted property, local
density---can be aggregated per node and rendered as a continuous background
field at a resolution set by $K$. Methods without a fixed latent grid can still
render a field by inverting an arbitrary raster, but have no canonical cells at
which to aggregate observed data. \emph{Heavy-tailed}: the low-dimensional
model has heavier-than-Gaussian tails, the usual remedy for crowding---built
into the similarity kernel of t-SNE (Student-t) and UMAP, and available in
\nfmap{} as a Student-t mixture prior. None of these methods give coordinates
that are comparable between separately trained models.}
\label{tab:capabilities}
\begin{tabular}{lccccc}
\toprule
 & Parametric & Invertible & Exact density & Landscape grid & Heavy-tailed \\
\midrule
t-SNE      & \ding{55} & \ding{55} & \ding{55} & \ding{55} & \ding{51} \\
UMAP       & \ding{51}$^{*}$ & \ding{55} & \ding{55} & \ding{55} & \ding{51} \\
GTM        & \ding{51} & \ding{55} & \ding{51} & \ding{51} & \ding{55} \\
FloDR      & \ding{51} & \ding{51} & \ding{51} & \ding{55}$^{\dagger}$ & \ding{55} \\
\nfmap{}   & \ding{51} & \ding{51} & \ding{51} & \ding{51} & \ding{51}$^{\ddagger}$ \\
\bottomrule
\end{tabular}

\vspace{2pt}
\footnotesize $^{*}$via the parametric variant; the standard implementation is
non-parametric. $^{\dagger}$FloDR's mixture components are learned rather than
fixed, so their positions and effective resolution vary with the data.
$^{\ddagger}$optional (\texttt{prior\_family='student\_t'}); the default prior is
a Gaussian mixture.
\end{table}
